\section{Introduction}
\label{sec:intro}

Breadth-first search (BFS) is one of the most basic graph algorithms. Starting from a source vertex, it visits the graph level by level, where level $d$ holds the vertices at distance $d$ from the source, and so finds shortest paths in unweighted graphs. It is a building block of many other graph algorithms, such as connectivity, betweenness centrality and maximum flow, and it is the search kernel of the Graph500 and GAP benchmarks~\cite{graph500, gapbs}.

In most parallel BFS algorithms for shared memory, the threads expand the vertices of one level together, which discovers the next level, and move on to it only once the whole level is done; such algorithms are \emph{level-synchronous}~\cite{bader_bfs, charles_bfs}. Beamer et al.~\cite{beamer_dobfs} observed that on low-diameter graphs, a few middle levels hold most of the vertices, and that on those levels it is cheaper to search \emph{bottom-up}: every undiscovered vertex looks for a neighbor in the current level and stops at the first one it finds, instead of every vertex of the level checking all of its neighbors. Their \emph{direction-optimizing} BFS chooses between the two directions at every level. It underlies the fastest shared-memory BFS implementations today, including those of GAPBS~\cite{gapbs}, Ligra~\cite{ligra}, GBBS~\cite{gbbs} and PASGAL~\cite{pasgal}.

These implementations update vertices with atomic instructions rather than locks, but their BFS methods are still \emph{blocking}: GAPBS ends every step with an OpenMP barrier, and GBBS and PASGAL end every round when its parallel loop joins, so a level is over only once its slowest thread has finished its share. A thread delayed while it holds part of a level, because the operating system preempted it, or an interrupt or a page fault stopped it, keeps every other thread waiting at the end of that level; dynamic scheduling and work stealing~\cite{blumofe_workstealing} hand over the work that a delayed thread has not started, but not the work it is in the middle of. The cost recurs at every level that a delay overlaps, much as short delays on single processors, amplified by synchronization, slow down whole applications on large parallel machines~\cite{petrini_noise, hoefler_noise}, and the number of levels grows with the diameter of the graph: a BFS of the USA road network in our experiments has 6,262 levels, and on such graphs the synchronization alone can cost more than the work~\cite{pasgal}.

Wait-freedom removes this dependence. A method is \emph{wait-free} if every call to it finishes in a finite number of the calling thread's own steps, however slow the other threads are~\cite{herlihy_waitfree, aomp}. Our BFS is a method that every thread calls with the same source, and that returns only once all distances are final; if it is wait-free, no call waits for other threads, and the search is complete as soon as the first call returns, even if the other threads have stopped for good.

A wait-free BFS is easy to obtain if work does not matter: every thread can search the whole graph by itself, all writing the same distances, at $T$ times the work of a sequential search with $T$ threads. Do-All and Certified Write-All (CWA), in which asynchronous or failing processors must complete a set of tasks~\cite{kanellakis_robust, martel_cwa, georgiou_doall}, show how to do far less, but the simplest way to use a CWA for BFS, one run per level over all $n$ vertex indices, costs at least $n$ in every level, however few vertices the level holds (Section~\ref{sec:bg-doall}). What has been missing is a wait-free BFS whose threads share each top-down level at a cost that depends on that level, and that is about as fast as the heavily tuned blocking implementations when no thread is delayed; wait-free implementations of concurrent data structures have long been considered too slow for practice~\cite{kogan_fastwf}.

Three problems make this hard. First, without barriers, every thread must decide for itself when a level is complete, and a thread that has fallen behind, possibly by several levels, can still write to shared state at any moment; such late writes must never undo or corrupt the work of others, and since the tasks of a level are discovered during the level before, by the same threads, there is no array of tasks to certify. Second, helping must be cheap: when no thread is delayed, threads should rarely repeat each other's work, and when one is, the others should split what it has left rather than each redo all of it, even when that is a vertex with millions of neighbors. Third, direction optimization adds decisions for a whole level: all threads must expand a level in the same direction, chosen from its size, which threads may see differently, and the bottom-up step needs a test of membership in the current level that stays exact for threads that are behind.

This paper presents a wait-free method for direction-optimizing BFS that solves these problems. Its contributions are:
\begin{itemize}
    \item \textbf{A wait-free direction-optimizing BFS} (Sections~\ref{sec:wf-topdown} to~\ref{sec:proof}, Figures~\ref{fig:td-protocol} and~\ref{fig:do-protocol}). Each level holds one bucket per thread; a thread expands its own bucket and then helps with the others in random orders, and every write that lets threads skip work comes after that work. One successful compare-and-swap (CAS) per level, with at most one attempt per thread, fixes its direction. The distances are final once any call returns, and every call returns within a bounded number of its own steps, under total store order and allocation in a bounded number of steps.
    \item \textbf{Bounds on work and memory} (Sections~\ref{sec:pf-work} and~\ref{sec:variants}, Figure~\ref{fig:mem-protocol}). The work of a top-down level depends only on that level and on $T$. Any piece of work certified as a whole can be repeated by every thread, so finer certificates can lower the repetition; with uniformly random orders and sequentially consistent or fenced marks, threads expand $O((T + R)\log n)$ positions of a bucket of $R$ positions with high probability, whatever the schedule. A variant that reuses two level nodes needs memory that depends only on the largest levels, plus a constant per level.
    \item \textbf{An implementation and evaluation} (Sections~\ref{sec:impl} and~\ref{sec:experiments}). On a 128-core machine, our method examines at most 30\% more neighbors than the sequential baseline at 112 threads on four real-world graphs and a Kronecker graph, and the best times of GAPBS, GBBS and PASGAL there are 1.58 to 2.77 times its own in the geometric mean. With one of 112 threads stalled for 10~ms at every level, or 1~ms on the road network, its median time grows by at most 7\%, and theirs up to 19.5-fold. Graphs with many levels of few vertices, such as long paths, are slow for it, and the reset before a search waits for every thread. Our code, the scripts of every experiment and their data are available at \codeurl.
\end{itemize}
To our knowledge, ours is the first wait-free BFS in which threads share the work of a search at a cost per top-down level that depends on the level rather than on the graph. On the social networks, web graph, road network, Kronecker graph and synthetic graph of hubs that we evaluate, its best times are below those of GAPBS, GBBS and PASGAL; on a long path, PASGAL's best is below that of our method.

Section~\ref{sec:background} gives background and states our model. Sections~\ref{sec:wf-topdown} and~\ref{sec:dir-opt} present the method, Section~\ref{sec:proof} proves it correct and wait-free and bounds its work, and Section~\ref{sec:variants} describes the variant that reuses level nodes to reduce space. Section~\ref{sec:impl} describes our implementation, Section~\ref{sec:experiments} evaluates it, Section~\ref{sec:related} discusses related work, and Section~\ref{sec:conclusion} concludes.
